% sections/abstract.tex -- abstract draft (includes the abstract environment; \input it before \maketitle)
\begin{abstract}
The renormalization-group functions $\beta$, $\eta$ and $\nu$ of the three-dimensional $O(N)$ model in fixed
dimension, in the massive zero-momentum scheme, were known to eighth order in the coupling. Here we extend them to
tenth order, and $\beta$ to eleventh, for general $N$---which requires the four-point and two-point functions to ten
loops, 1\,593\,254 graphs at the last order---and we re-evaluate all lower orders independently. Through the
Callan--Symanzik equation every coefficient follows algebraically from per-order sums that are linear in finite,
zero-momentum $\Rbar$-subtracted parametric integrals, one per graph; we numerically sample them on GPUs by tropical Monte Carlo, with densities adapted to the
forest terms. The computation reproduces the known per-order sums from two to seven loops within three standard
deviations, and agrees at ninth and tenth order, at lower precision, with an independent route through the vertex
functions of an intermediate scheme. The new
coefficients have relative uncertainties of a few $10^{-5}$ or better, and the eighth-order coefficients of $\eta$ and
$\nu$ improve by factors of \etaEightGain\ and \nuEightGain. An encouraging result concerns the cost: with this strategy,
the time needed to integrate an order to a fixed relative precision grows from seven loops on by a factor of only
\costgrowthrange\ per loop on average. However, the construction of the integrands still grows factorially, and at tenth order it costs more than the
integration: \costtenthcpu\ CPU core-hours against \costtenthgpu\ GPU hours. Resummed, the tenth-order series give
$\nu=\rsNuDNone$, $\eta=\rsEtaDNone$ and $\omega=\rsOmegaDNone$ for $N=1$, against $\rsBenchCBNuNone$,
$\rsBenchCBEtaNone$ and $\rsBenchCBOmegaNone$ from the conformal bootstrap; for $N=2$ and 3 the deviation of $\nu$ grows, to
\rsPullNuNthree\ standard deviations at $N=3$, where $\omega$ lies \rsPullOmegaNthree\ standard deviations from the
Monte Carlo value, so that the uncertainty is now dominated by the resummation
rather than by the coefficients. Our computational strategy is cooperative
and distributed, and thus open to future improvement: every result presented here is generated automatically from, or
checked against, a public ledger of independent, exactly pooled runs, hosted on GitHub, to which anyone is welcome to contribute
computing time.
\end{abstract}
